\section{Заключение}

Разработана и экспериментально проверена воспроизводимая схема послойного CPU-инференса GPT-2-совместимой модели с использованием Ray Actors. Реализованы четыре варианта обработки выходных логитов D0--D3 и два монолитных контроля B0/B1; выполнено 2160 сравнительных измерений основной матрицы нагрузок и 1296 дополнительных измерений монолитного контроля. Во всех этих запусках получено совпадение генерируемых токенов с эталоном.

Выделен существенный эффект исключения лишних логитов из вычисления и передачи. Однако показано, что для очень маленькой модели на единственной физической машине оптимизированное распределённое исполнение остаётся медленнее равноценно оптимизированного монолитного. Следовательно, сокращение полезного объёма передаваемых данных само по себе не гарантирует выигрыша по полной задержке.

Математическая постановка выбора границ разбиения с ограничениями памяти реализована методом динамического программирования, проверена полным перебором и применена к работе шестислойной DistilGPT2. Получены корректные результаты распределённой генерации для двух вариантов размещения. Проверка преимуществ на реальных неоднородных серверах остаётся открытой задачей: используемый профиль стоимости был синтетическим, а измерения выполнялись на одной машине.

Следующими этапами являются получение реальных профилей памяти и вычислительных затрат, повторение эксперимента на доступном MacBook и при возможности на независимых физических узлах, внедрение KV-кэша и исследование многозапросного конвейера. Реализованная экспериментальная инфраструктура и авторский набор сценариев предоставляют воспроизводимую основу для этих исследований.

